\documentclass[12pt,a4paper]{article}

% --- Настройка языка и шрифтов для LuaLaTeX ---
\usepackage{fontspec}
\defaultfontfeatures{Ligatures=TeX,Scale=MatchLowercase}

\setmainfont{Alice-Regular}[
    Path=font/,
    Extension=.ttf,
    BoldFont = Alice-Regular,
    ItalicFont = Alice-Regular,
    BoldItalicFont = Alice-Regular,
    BoldFeatures = {FakeBold=1.6},
    ItalicFeatures = {FakeSlant=0.2},
    BoldItalicFeatures = {FakeBold=1.6,FakeSlant=0.2}
]

\newfontfamily\cyrillicfont{Alice-Regular}[
    Path=font/,
    Extension=.ttf,
    BoldFont = Alice-Regular,
    ItalicFont = Alice-Regular,
    BoldItalicFont = Alice-Regular,
    BoldFeatures = {FakeBold=1.6},
    ItalicFeatures = {FakeSlant=0.2},
    BoldItalicFeatures = {FakeBold=1.6,FakeSlant=0.2}
]

\usepackage{polyglossia}
\setmainlanguage{russian}
\setsansfont{TeX Gyre Heros}
\setmonofont{TeX Gyre Cursor}

\usepackage{graphicx}        					% Графика
\usepackage{amsmath, amsfonts, amssymb, amsthm}	% Математика
\usepackage{mathtools}
\usepackage{chemfig}
\usepackage[version=4,arrows=pgf-filled,
            textfontname=sffamily,
            mathfontname=mathsf]{mhchem}
\usepackage[protrusion=false]{microtype}       % Микротипографика
\usepackage{braket} 						% Braket нотация для квантовой механики
\usepackage{unicode-math} 					% Современный пакет для математики в LuaLaTeX
\usepackage{longtable}       				% Таблицы с разрывом страниц
\usepackage{tikz}           				% Рисование
\usepackage[europeanresistors]{circuitikz}	% Рисование электрических схем в европейском стиле
\usepackage{listings}						% Для отображения кода
\usepackage{tabularx}       				% Для таблиц во всю ширину
\usepackage{array}		  					% Для расширенных возможностей таблиц
\usepackage{float}          				% Для точного позиционирования фигур
\usepackage{pgfplots}                          % Для построения графиков
\usepackage{caption}                           % Для настройки подписей к рисункам и таблицам
\usepackage{xfp}                               % Для вычислений в LaTeX
\usepackage{luacode}                           % Для выполнения Lua-кода в документе    
\pgfplotsset{compat=1.18}
\usepackage[paper=letterpaper,
            top=20mm,
            bottom=20mm,
            left=10mm,
            right=10mm,
            footskip=0.4in 
]{geometry}

\usetikzlibrary{arrows.meta}

% --- Колонтитулы ---
\usepackage{fancyhdr}
\pagestyle{fancy}
\fancyhf{}

% --- Теоремы и окружения ---
\theoremstyle{plain}
\newtheorem{theorem}{Theorem}[section]
\newtheorem{lemma}[theorem]{Lemma}
\newtheorem{proposition}[theorem]{Proposition}
\newtheorem{corollary}[theorem]{Corollary}

\theoremstyle{definition}
\newtheorem{definition}[theorem]{Definition}

\theoremstyle{remark}
\newtheorem{remark}[theorem]{Remark}

\numberwithin{equation}{section}


\renewcommand{\headrulewidth}{0.4pt}
\renewcommand{\footrulewidth}{0.4pt}
\renewcommand\lstlistingname{Исходный код}

\lstdefinestyle{main}{
	backgroundcolor=\color{gray!8},   
	commentstyle=\color{green!50!black},
	keywordstyle=\color{blue!80!black},
	numberstyle=\small\color{darkgray},
	stringstyle=\rmfamily\color{orange!90!black},
	basicstyle=\ttfamily\small,
	breakatwhitespace=false,         
	breaklines=true,                   
	keepspaces=true,                 
	numbers=left,                    
	numbersep=8pt,                  
	showspaces=false,                
	showstringspaces=false,
	showtabs=false,                  
	tabsize=4
}	

\lstset{style=main}

\newcolumntype{C}{>{\centering\arraybackslash}X}
\renewcommand{\arraystretch}{1.25}

\begin{luacode*}
    data = require("data")
    include = require("include")
    handle = require("handle")
\end{luacode*}

\begin{document}

\section{\large\textnormal{В объеме полупроводникового прибора генерируются 
электронно-дырочные пары. В момент времени $t = 0$ генерация прекращается. \\
- В момент времени $t_1$ концентрация избыточных носителей $\Delta n(t_1)$, \\
- В момент времени $t_2$ концентрация избыточных носителей $\Delta n(t_2)$. \\
Избыточные носители распределены равномерно по объему полупроводника. \\
Значения $t_1$, $t_2$, $\Delta n(t_1)$, $\Delta n(t_2)$ заданы в соответствии с вариантом. \\
Найти: $\Delta n(0)$ и $\tau_n$ (время жизни носителей заряда – электронов)}}

\begin{figure}[H]
\centering
\begin{minipage}[t]{0.3\textwidth}
    \vspace{0em}
    \begin{tabularx}{\textwidth}{|X|}
        \hline
        \multicolumn{1}{|C|}{Дано} \\
        \hline
        $t_1 = 10^{-7}$ с \\
        $t_2 = 10^{-6}$ с \\
        $\Delta n(t_1) = 10^{18}$ см$^{-3}$ \\
        $\Delta n(t_2) = 10^{17}$ см$^{-3}$ \\
        $t = 0$ с — генерация прекращается \\
        \hline
        \multicolumn{1}{|C|}{Найти} \\
        \hline
        $\Delta n(0)$ — ? \\
        $\tau_n$ — ? \\
        \hline
    \end{tabularx}

    \begin{equation*}
        \begin{cases}
            \Delta n(t_1) = \Delta n(0) e^{-\frac{t_1}{\tau_n}} \\
            \Delta n(t_2) = \Delta n(0) e^{-\frac{t_2}{\tau_n}}
        \end{cases}
    \end{equation*}
    
    $\tau_n$ легко выразить как
    \begin{equation*}
        \begin{aligned}
            \frac{\Delta n(t_1)}{\Delta n(t_2)} = e^{\frac{t_2 - t_1}{\tau_n}} \Longrightarrow 
            \\ 
            \Longrightarrow \tau_n = \frac{t_2 - t_1}{\ln\left(\frac{\Delta n(t_1)}{\Delta n(t_2)}\right)}
        \end{aligned}
    \end{equation*}
    из известного $\tau_n$ можно найти
    \begin{equation*}
        \Delta n(0) = \Delta n(t_1) e^{\frac{t_1}{\tau_n}} = \Delta n(t_2) e^{\frac{t_2}{\tau_n}}
    \end{equation*}

    \begin{equation*}
        \tau_n = \frac{10^{-6} - 10^{-7}}{\ln\left(\frac{10^{18}}{10^{17}}\right)} = 
        \directlua{
            tau_n = (10^(-6) - 10^(-7)) / math.log(10^(18) / 10^(17))
        
            tex.sprint(include.value_with_mantice(tau_n, 2, 0, 0))
        }
        \text{ с}
    \end{equation*}
    \begin{equation*}
        \displaystyle{\Delta n(0) = 10^{18} e^{\frac{10^{-7}}{\directlua{
            tex.sprint(include.value_with_mantice(tau_n, 2, 0, 0))
        }}} = \directlua{
            delta_n_0 = 10^(18) * math.exp(10^(-7) / tau_n)

            tex.sprint(include.value_with_mantice(delta_n_0, 2, 0, 0))
        }
        \text{ см$^{-3}$}}
    \end{equation*}
\end{minipage}\hfill\begin{minipage}[t]{0.7\textwidth}
\vspace{0em}
\begin{tabularx}{\textwidth}{C|}
    \hline
    Решение \\
    \hline
\end{tabularx}

\vspace{0.5em}
\hspace{0.5em}
\begin{minipage}[t]{0.95\textwidth}
$n_0$ — концентрация носителей заряда в состоянии равновесия (в отсутствии генерации) 

\begin{equation*}
    \frac{d(\Delta n)}{dt} = -\frac{\Delta n}{\tau_n} \Longrightarrow \Delta n(t) = \Delta n(0) e^{-\frac{t}{\tau_n}}
\end{equation*}

\directlua{
    local func = handle.delta_n:new(delta_n_0, tau_n)
    local x = { 0.1e-7, 1e-7, 5e-7, 1e-6, 2e-6, 5e-6 }
    local table_delta_n = {
        "red!70!black",
        "*",
        "2pt",
        nil,
        x, func:get_values(x)
    }

    include.pgf_plot_wrap("График зависимости концентрации избыточных носителей от времени", 
        "Время, с", "Концентрация избыточных носителей, см$^{-3}$", "north east", 1, 1, table_delta_n)
    
}

\end{minipage}\end{minipage}

\end{figure}

Ответ: $\Delta n(0) = \directlua{
    tex.sprint(include.value_with_mantice(delta_n_0, 2, 0, 0))
}$ см$^{-3}$, $\tau_n = \directlua{
    tex.sprint(include.value_with_mantice(tau_n, 2, 0, 0))
}$ с

\pagebreak

\section{\large\textnormal{В соответствии с предыдущим заданием построить график зависимости
$\frac{\Delta n}{\Delta n(0)}$ от $\frac{t}{\tau_n}$}}

\directlua{
    local func = handle.delta_n:new(delta_n_0, tau_n)
    local x_norm = {}
    local x = { 1, 2, 3, 4, 5, 6, 7, 8, 9, 10 }
    for _, v in ipairs(x) do
        table.insert(x_norm, v * tau_n)
    end
    local y = func:get_values(x_norm)
    for i, v in ipairs(y) do
        y[i] = v / delta_n_0
    end
    local table_delta_n = {
        "blue!70!black",
        "*",
        "2pt",
        nil,
        x, y
    }

    include.pgf_plot_wrap("График зависимости $\\frac{\\Delta n}{\\Delta n(0)}\\left(\\frac{t}{\\tau_n}\\right)$", 
        "$\\frac{t}{\\tau_n}$", "$\\frac{\\Delta n}{\\Delta n(0)}$", "north east", 0.9, 1.1, table_delta_n)
}

\pagebreak

\section{\large\textnormal{Полупроводник \protect\chemfig{n-Si} освещается, заданы в соответствии с вариантом:\\
- Скорость генерации электронно-дырочных пар $G$, \\
- Время жизни электронов $\tau_n$, \\
- Равновесная концентрация электронов $n_0$, \\
- $T = 300$ К. \\
Найти: $\frac{\Delta n}{n_0}$ и $\frac{\Delta p}{p_0}$ в установившемся режиме}}

\begin{figure}[H]
    \begin{minipage}[t]{0.3\textwidth}
        \vspace{0em}
        \begin{tabularx}{\textwidth}{|X|}
            \hline
            \multicolumn{1}{|C|}{Дано} \\
            \hline
            $G = 10^{19}$ см$^{-3}$ с$^{-1}$ \\
            $\tau_n = 10^{-7}$ с \\
            $n_0 = 10^{18}$ см$^{-3}$ \\
            $T = 300$ К \\
            $n_i = 1.45 \cdot 10^{10}$ см$^{-3}$ \\
            \hline
            \multicolumn{1}{|C|}{Найти} \\
            \hline
            $\frac{\Delta n}{n_0}$ — ? \\
            $\frac{\Delta p}{p_0}$ — ? \\
            \hline
        \end{tabularx}

        \vspace{0.5em}
        
        В установившемся режиме:
        \begin{equation*}
            G = R_n = R_p
        \end{equation*}
        где $R_n$ и $R_p$ — скорости рекомбинации электронов и дырок соответственно.

        Тогда
        \begin{equation*}
            \begin{aligned}
                R_n = \frac{\Delta n}{\tau_n} \Longrightarrow \Delta n = G \tau_n \\
                R_p = \frac{\Delta p}{\tau_n} \Longrightarrow \Delta p = G \tau_n
            \end{aligned}
        \end{equation*}
    \end{minipage}\hfill\begin{minipage}[t]{0.7\textwidth}
        \vspace{0em}
        \begin{tabularx}{\textwidth}{C|}
            \hline
            Решение \\
            \hline
        \end{tabularx}

        \vspace{0.5em}
        \hspace{0.5em}
        \begin{minipage}[t]{0.95\textwidth}
            $n_i$ — собственная концентрация носителей заряда в 
            полупроводнике \chemfig{Si} при температуре $T = 300$ К.
            
            Воспользуемся законом действующих масс:
            \begin{equation*}
                n_0 p_0 = n_i^2 \Longrightarrow p_0 = \frac{n_i^2}{n_0} = \directlua{
                    n_0 = 10^(18)
                    n_i = 1.45 * 10^(10)
                
                    p_0 = (n_i)^2 / n_0

                    tex.sprint("\\frac{(" 
                               .. include.value_with_mantice(n_i, 2, 0, 0) .. 
                               ")^2}{" 
                               .. include.value_with_mantice(n_0, 2, 0, 0) .. 
                               "} = " 
                               .. include.value_with_mantice(p_0, 2, 2, 0))
                } \text{ см$^{-3}$}
            \end{equation*}

            Так как режим установившийся, то
            \begin{equation*}
                \Delta n = \Delta p = G \tau_n = \directlua{
                    G = 10^(19)
                    tau_n = 10^(-7)
                
                    delta_n = G * tau_n
                
                    tex.sprint(include.value_with_mantice(delta_n, 2, 0, 0))
                } \text{ см$^{-3}$} 
            \end{equation*}

            Тогда
            \begin{equation*}
                \begin{aligned}
                    \frac{\Delta n}{n_0} = \frac{G \tau_n}{n_0} = \directlua{
                        d_delta_n = delta_n / n_0

                        tex.sprint("\\frac{" 
                                   .. include.value_with_mantice(delta_n, 2, 0, 0) .. 
                                   "}{" 
                                   .. include.value_with_mantice(n_0, 2, 0, 0) .. 
                                   "} = " 
                                   .. include.value_with_mantice(d_delta_n, 2, 0, 0))
                    } \text{ см$^{-3}$} \\
                    \frac{\Delta p}{p_0} = \frac{G \tau_n}{p_0} = \directlua{
                        d_delta_p = delta_n / p_0

                        tex.sprint("\\frac{" 
                                   .. include.value_with_mantice(delta_n, 2, 0, 0) .. 
                                   "}{" 
                                   .. include.value_with_mantice(p_0, 2, 2, 0) .. 
                                   "} = " 
                                   .. include.value_with_mantice(d_delta_p, 2, 0, 0))
                    } \text{ см$^{-3}$}
                    \end{aligned}
            \end{equation*}
    
            \vspace{2em}
            Ответ: $\frac{\Delta n}{n_0} = \directlua{
                tex.sprint(include.value_with_mantice(d_delta_n, 2, 0, 0))
            }$ и $\frac{\Delta p}{p_0} = \directlua{
                tex.sprint(include.value_with_mantice(d_delta_p, 2, 0, 0))
            }$.
    \end{minipage}\end{minipage}
\end{figure}

\pagebreak

\section{\large\textnormal{Дан параллелепипед из \protect\chemfig{Si}. В соответствии 
с вариантом заданы: \\
- Тип полупроводника, \\
- Концентрация примеси, \\
- Длина и площадь поперечного сечения полупроводника. \\
Найти: удельное сопротивление $\rho$, полное сопротивление $R$.}}

\begin{figure}[H]
    \begin{minipage}[t]{0.3\textwidth}
        \vspace{0em}
        \begin{tabularx}{\textwidth}{|X|}
            \hline
            \multicolumn{1}{|C|}{Дано} \\
            \hline
            Полупроводник $n$-типа \\
            $N_D = 10^{20}$ см$^{-3}$ \\
            $\frac{L}{S} = 10^{-3}$ мкм$^{-1}$ \\
            \hline
            \multicolumn{1}{|C|}{Найти} \\
            \hline
            $\rho$ — ? \\
            $R$ — ? \\
            \hline
        \end{tabularx}

        \directlua{
            N_D = 10^(20)
            L_S = 10 
            q = 1.6 * 10^(-19)
        }

        \vspace{0.5em}
        Полное сопротивление проводника определяется как
        \begin{equation*}
            R = \frac{\rho L}{S}
        \end{equation*}

        Удельное сопротивление определяется как
        \begin{equation*}
            \rho = \frac{1}{q n_0 \mu_n}
        \end{equation*}
        
    \end{minipage}\hfill
    \begin{minipage}[t]{0.7\textwidth}
        \vspace{0em}
        \begin{tabularx}{\textwidth}{C|}
            \hline
            Решение \\
            \hline            
        \end{tabularx}

        \vspace{0.5em}
        \hspace{0.025\textwidth}
        \begin{minipage}[t]{0.95\textwidth}
            \indent При комнатной температуре $T = 300$ К мелкая донорная примесь полностью ионизирована,
            поэтому концентрация электронов в полупроводнике $n$-типа равна концентрации донорной примеси
            \begin{equation*}
                n_0 = N_D = 10^{20} \text{ см$^{-3}$}
            \end{equation*}
            Подвижность электронов в кремнии при сильной легированности примеси $N_D = 10^{20}$ см$^{-3}$ 
            и соотвественно преобладанию рассеяния на примеси, равна
            \begin{equation*}
                \mu_n = 80 \text{ см$^2$/(В·с)}
            \end{equation*}
            Соответственно, удельное сопротивление 
            \begin{equation*}
                \rho = \frac{1}{q n_0 \mu_n} = \directlua{
                    mu_n = 80

                    rho = 1 / (q * N_D * mu_n)

                    tex.sprint("\\frac{1}{(" 
                               .. include.value_with_mantice(q, 2, 0, 0) .. 
                               ") \\cdot (" 
                               .. include.value_with_mantice(N_D, 2, 0, 0) .. 
                               ") \\cdot (" 
                               .. include.value_with_mantice(mu_n, 2, 1, 0) .. 
                               ")} = " 
                               .. include.value_with_mantice(rho, 2, 0, 0))
                } \text{ Ом·см}
            \end{equation*}

            Полное сопротивление проводника
            \begin{equation*}
                R = \frac{\rho L}{S} = \directlua{
                    R = rho * L_S

                    tex.sprint(include.value_with_mantice(rho, 2, 0, 0) .. " \\cdot " 
                               .. include.value_with_mantice(L_S, 2, 1, 0) .. 
                               " = " 
                               .. include.value_with_mantice(R, 2, 0, 0))
                } \text{ Ом}
            \end{equation*}
        \end{minipage}
    \end{minipage}

    Ответ: $\rho = \directlua{
        tex.sprint(include.value_with_mantice(rho, 2, 0, 0))
    }$ Ом·см и $R = \directlua{
        tex.sprint(include.value_with_mantice(R, 2, 0, 0))
    }$ Ом
\end{figure}

\pagebreak

\section{\large\textnormal{Для \protect\chemfig{Si}, $n$-типа ($T=300К$), длиной $L$ (см. вариант), 
заданы распределения примеси: \\
а) $N_{D1}(x) = N_{Ds}$, \\
б) $N_{D2}(x) = N_{Ds} \cdot e^{-x/L}$, \\
в) $N_{D3}(x) = N_{Ds} \cdot e^{-x^2/L^2}$. \\
Найти напряжённость электрического поля $E(x=L)$ для случаев а-в.}}

\begin{figure}[H]
    \begin{minipage}[t]{0.3\textwidth}
        \vspace{0em}
        \begin{tabularx}{\textwidth}{|X|}
            \hline
            \multicolumn{1}{|C|}{Дано} \\
            \hline
            Полупроводник $n$-типа \\
            $T = 300$ К \\
            $L = 10$ мкм \\
            \hline
            \multicolumn{1}{|C|}{Найти} \\
            \hline
            $E(x=L)$ — ? \\
            \hline
        \end{tabularx}

        \vspace{0.5em}
        Формлула напряжённости электрического поля в полупроводнике с градиентом концентрации примеси
        \begin{equation*}
            E(x) = -\frac{kT}{q} \cdot \frac{1}{N_D(x)} \cdot 
                \frac{dN_D(x)}{dx}
        \end{equation*}
        
    \end{minipage}\hfill
    \begin{minipage}[t]{0.7\textwidth}
        \vspace{0em}
        \begin{tabularx}{\textwidth}{C|}
            \hline
            Решение \\
            \hline            
        \end{tabularx}

        \vspace{0.5em}
        \hspace{0.025\textwidth}
        \begin{minipage}[t]{0.95\textwidth}
            \indent Для случая а) $N_{D1}(x) = N_{Ds}$, тогда
            \begin{equation*}
                E(x=L) = -\frac{kT}{q} \cdot \frac{1}{N_{Ds}} \cdot 0 = 0
            \end{equation*}

            Для случая б) $N_{D2}(x) = N_{Ds} \cdot e^{-x/L}$, тогда
            \begin{equation*}
                \begin{aligned}
                    E(x=L) &= -\frac{kT}{q} \cdot \frac{1}{N_{Ds} e^{-L/L}} \cdot 
                        (-\frac{N_{Ds}}{L} e^{-L/L}) = -\frac{kT}{qL} \\
                    E(x=L) &= \directlua{
                        k = 1.38 * 10^(-23)
                        T = 300
                        q = 1.6 * 10^(-19)
                        L = 10 * 10^(-6)

                        E = -k * T / (q * L)

                        tex.sprint("-\\frac{" 
                                   .. include.value_with_mantice(k, 2, 0, 0) ..
                                   "\\cdot" 
                                   .. include.value_with_mantice(T, 2, 0, 0) .. 
                                   "}{"
                                   .. include.value_with_mantice(q, 2, 0, 0) ..
                                   "\\cdot" 
                                   .. include.value_with_mantice(L, 2, 0, 0) .. 
                                   "} = "
                                   .. include.value_with_mantice(E, 2, 0, 0) .. 
                                   "\\text{ В/м}")
                    }
                \end{aligned}
            \end{equation*}

            Для случая в) $N_{D3}(x) = N_{Ds} \cdot e^{-x^2/L^2}$, тогда
            \begin{equation*}
                \begin{aligned}    
                E(x=L) &= -\frac{kT}{q} \cdot \frac{1}{N_{Ds} e^{-L^2/L^2}} \cdot 
                    (-\frac{2L}{L^2} N_{Ds} e^{-L^2/L^2}) = -\frac{2kT}{qL} \\
                E(x=L) &= \directlua{
                        k = 1.38 * 10^(-23)
                        T = 300
                        q = 1.6 * 10^(-19)
                        L = 10 * 10^(-6)

                        E = -2 * k * T / (q * L)

                        tex.sprint("-\\frac{2 \\cdot " 
                                   .. include.value_with_mantice(k, 2, 0, 0) ..
                                   "\\cdot" 
                                   .. include.value_with_mantice(T, 2, 0, 0) .. 
                                   "}{"
                                   .. include.value_with_mantice(q, 2, 0, 0) ..
                                   "\\cdot" 
                                   .. include.value_with_mantice(L, 2, 0, 0) .. 
                                   "} = "
                                   .. include.value_with_mantice(E, 2, 0, 0) .. 
                                   "\\text{ В/м}")
                }
                \end{aligned}
            \end{equation*}
    \end{minipage}\end{minipage}

    \vspace{1em}
    Ответ: для случая а $E(x=L) = 0$ В/м, для случая б $E(x=L) = \directlua{
        tex.sprint(include.value_with_mantice(E, 2, 0, 0))
    }$ В/м, для случая в $E(x=L) = \directlua{
        tex.sprint(include.value_with_mantice(E, 2, 0, 0))
    }$ В/м.
\end{figure}

\end{document}